In this section we review necessary background for the proof of our main theorem---Theorem~\ref{thm: d4 hj3}. We begin by recalling the Miller square, introduced in \cite{MillerSquare}, which captures the interplay between the Adams spectral sequence and the Adams--Novikov spectral sequence. We then discuss the recent cofiber of $tau$ method developed by Gheorghe, Isaksen, Wang and the second author \cite{GWX, IWX, IWX2}, which uses the motivic stable homotopy category over $\mathbb{C}$ and the motivic Adams spectral sequence to categorify the Miller square. %Reader who are familiar with both material should feel free to skip this section.

\subsection{The Miller square}\hfill

The Adams spectral sequence and the Adams--Novikov spectral sequence are two of the most effective methods of computing the homotopy groups of the $p$-completed sphere spectrum, $S^0$. They are spectral sequences of the form:
\begin{align*}
  \Ext_{\A}^{s,t}(\F_p,\, \F_p) &\cong E_2^{s,t} \Longrightarrow \pi_{t-s}S^0, \ &&  d_r:E_r^{s,t} \rightarrow E_r^{s+r, t+r-1} \\
  \Ext_{\BP_*\BP}^{s,t}(\BP_*,\, \BP_*) &\cong E_2^{s,t} \Longrightarrow \pi_{t-s}S^0, \quad &&  d_r:E_r^{s,t} \rightarrow E_r^{s+r, t+r-1}
\end{align*}
where $\A$ is the mod $p$ dual Steenrod algebra and $\BP$ is the Brown--Peterson spectrum at the prime $p$. For degrees, $s$ is the homological degree and is referred as the Adams filtration (resp. the Adams--Novikov filtration), and $t$ is the internal degree.


It is important to understand connections between them. A first connection is given by the Thom reduction map $\BP\rightarrow \F_p$, which induces a map of spectral sequences
$$\Ext_{\BP_*\BP}^{s,t}(\BP_*,\, \BP_*) \longrightarrow \Ext_{\A}^{s,t}(\F_p,\, \F_p)$$
that preserves the $(s,t)$-degrees.
However, a general homotopy class in $\pi_{*}S^0$ doesn't usually have the same Adams filtration as the Adams--Novikov filtration. So this map is not very useful for comparison of the Adams filtration and the Adams--Novikov filtration of a surviving homotopy class---it only tells us the latter is less or equal to the former.

A fundamental connection is the Miller square. We have an algebraic Novikov spectral sequence converging to the Adams--Novikov $E_2$-page, and a Cartan--Eilenberg spectral sequence converging to the Adams $E_2$-page. It turns out the $E_2$-pages of these two algebraic spectral sequences are isomorphic.

The algebraic Novikov spectral sequence comes the filtration of powers of the augmentation ideal $I = (p, v_1, v_2, \cdots) \subset \BP_*$. It has the form:
$$ \Ext_{\BP_*\BP/I}^{s,t'}(\BP_*/I,\, I^k/I^{k+1}) \cong E_2^{s,k, t'} \Longrightarrow \Ext_{\BP_*\BP}^{s,t'}(\BP_*,\, \BP_*)$$
$$d_r: E_r^{s,k,t'} \longrightarrow E_r^{s+1,k+r-1,t'}$$
where $s$ and $k$ are homological degrees, and $t'$ is internal degree. In particular, in the Adams--Novikov gradings, all differentials in the algebraic Novikov spectral sequence look like Adams--Novikov $d_1$-differentials.

For $p=2$, let $\P$ be the sub-Hopf algebra of squares inside $\A$
\[ \P \cong \F_2[\xi_1^2, \xi_2^2, \cdots] \subseteq \F_2[\xi_1, \xi_2, \cdots] \cong \A \]
and let $\mathcal{Q}$ be the quotient Hopf algebra,
$\mathcal{Q} \cong \A \otimes_{\P} \F_2 \cong \Lambda_{\F_2}[\xi_1, \xi_2, \cdots]$.
The Cartan-Eilenberg spectral sequence comes from the extension of Hopf algebras
$\P \to \A \to \mathcal{Q}$.
It has the form:
$$E_2^{s,k,t} = \Ext^{s,t}_{\P}(\F_2, \Ext_{\mathcal{Q}}^k(\F_2, \F_2)) \Longrightarrow \Ext_{\A}^{s+k, t}(\F_2, \F_2)$$
$$d_r: E_r^{s,k,t} \longrightarrow E_r^{s+r,k-r-1,t}$$
where $s$ and $k$ are homological degrees, and $t$ is internal degree. In particular, in the Adams gradings, all differentials in the Cartan--Eilenberg spectral sequence look like Adams $d_1$-differentials.
The Hopf algebra $\mathcal{Q}$ is cocommutative and primitively generated by the exterior classes $\xi_{i+1}$, therefore we have
$$\Ext_{\mathcal{Q}}^{*,*}(\F_2, \F_2) \cong \F_2[q_0, q_1, \cdots],$$
where $q_i$ corresponds to $[\xi_{i+1}]$ and has $(k,t)$ bidegree $(1, 2^{i+1}-1)$.

We identify the $E_2$-pages of the Cartan--Eilenberg spectral sequence and the algebraic Novikov spectral sequence by using the isomorphism of Hopf algebroids
\[ (\BP_*/I,\ \BP_*\BP/I) \cong (\F_2,\ \P) \]
and the associated isomorphism of $\P$-comodule algebras
\[ \Ext_{\mathcal{Q}}^*(\F_2, \F_2) \cong \F_2[q_0, q_1, \cdots] \cong \F_2[v_0, v_1, \cdots] \cong \oplus_* I^*/I^{*+1} \]
%% \begin{align*}
%% \BP_*/I & \cong \F_p,\\
%% \BP_*\BP/I & \cong P,\\
%% \Ext_Q^*(\F_p, \F_p) & \cong \F_p[q_0, q_1, \cdots],\\
%% I^*/I^{*+1} & \cong \F_p[v_0, v_1, \cdots].
%% \end{align*}
where the middle isomorphisms identifies $q_i$ with $v_i$.
Taking into account that $q_i$ has $(s,k,t)$-degree $(0,1,2^{i+1} - 1)$ and $v_i$ has $(s,k,t')$-degree $(0,1,2^{i+1} - 2)$ the isomorphisms above provide an isomorphism
$$\xymatrix{
\Ext^{s,t}_{\P}(\F_2, \Ext_{\mathcal{Q}}^k(\F_2, \F_2)) \ar[rr]^-{\cong} & & \Ext_{\BP_*\BP/I}^{s,t'}(\BP_*/I, I^k/I^{k+1})
}$$
by sending $s$ to $s$, $k$ to $k$, and $t$ to $t'+k$.
Altogether, we have introduced the Miller square:

%% We further identify $q_i$ with $v_i$ keeping in mind that their $s$ and $k$-degrees are 0 and 1, and their $t$ and $t'$-degrees are $2^{i+1} - 1$ and $2^{i+1} - 2$. Then we have an isomorphism
%% $$\xymatrix{
%% \Ext^{s,t}_P(\F_p, \Ext_Q^k(\F_p, \F_p)) \ar[rr]^-{\cong} & & \Ext_{\BP_*\BP/I}^{s,t'}(\BP_*/I, I^k/I^{k+1})
%% }$$
%% by sending $s$ to $s$, $k$ to $k$, and $t$ to $t'+k$. Therefore, we have the Miller square.

\begin{displaymath}
    \xymatrix{
  & & \Ext^{s,t}_{\P}(\F_2, \Ext_{\mathcal{Q}}^k(\F_2, \F_2)) \ar@{=>}[ddll]|{\mathbf{Cartan\textup{-}Eilenberg \ SS}} \ar@{=>}[ddrr]|{\mathbf{Algebraic \ Novikov \ SS}} & &  \\
  & & & & \\
 \Ext_{\A}^{s+k,t}(\F_2, \F_2) \ar@{=>}[ddrr]|{\mathbf{Adams \ SS}} & & & & \Ext^{s,t-k}_{\BP_*\BP}(\BP_*,\BP_*) \ar@{=>}[ddll]|{\mathbf{Adams\textup{-}Novikov \ SS}}\\
  & & & & \\
  & & \pi_{t-s-k}S^0  & &
    }
\end{displaymath}

\begin{rmk}
  Miller's original motivation for studying his square was in order to deduce $d_2$-differentials in the Adams spectral sequence from $d_2$-differentials in the algebraic Novikov spectral sequence (see \cite{MillerSquare}). Using motivic homotopy theory, one can generalize this method to obtain information about $d_r$-differentials for $r \geq 2$. We will explain this in the next subsection.
\end{rmk}

\begin{rmk}
  At odd primes, the dual Steenrod algebra admits an additional \emph{Cartan grading} which places $\xi_i$ in degree $0$ and $\tau_i$ in degree $1$. As the differentials in the Cartan--Eilenberg spectral sequence must respect the Cartan grading, this spectral sequence collapses at the $E_2$-page (see \cite{RavenelGreenBook} for example).
\end{rmk}


\subsection{Motivic homotopy theory} \label{section2.2}   \hfill

We work with the motivic stable homotopy category over $\mathbb{C}$. For the gradings, we denote by $S^{1,0}$ the simplicial sphere, and by $S^{1,1}$ the multiplicative group $\mathbb{G}_m = \mathbb{A}^1 - 0$. We use the same notation for their suspension spectra. We denote by $\F_p^{\textup{mot}}$ the mod $p$ motivic Eilenberg-Mac Lane spectrum that represents the mod $p$ motivic cohomology, and by $\textup{BPGL}$ the motivic Brown-Peterson spectrum at the prime $p$. When the context is clear, we abuse notation and also write $S^{n,w}$ for the $\F_p^{\textup{mot}}$-completed motivic sphere spectrum in bidegree $(n,w)$. 

There is a map
$$\tau: S^{0,-1} \longrightarrow S^{0,0}$$
that induces a nonzero map on mod $p$ motivic homology (\cite{Voe032}, \cite{HuKrizOrmsby}), which has played a significant role in the $\mathbb{C}$-motivic stable homotopy theory \cite{Isaksen}, \cite{GWX} \cite{IWX}. We denote by $S^{0,0}/\tau$ the cofiber of $\tau$.
$$\xymatrix{   
S^{0,-1} \ar[r]^\tau & S^{0,0} \ar[r] & S^{0,0}/\tau \ar[r] & S^{1,-1}.  
}$$

There is a Betti Realization functor $\mathbf{Re}$ from the motivic stable homotopy category over $\mathbb{C}$ to the classical stable homotopy category. We have 
\begin{align*}
\mathbf{Re}(S^{n,w}) & \simeq S^n,\\
\mathbf{Re}(\F_p^\textup{mot}) & \simeq \F_p,\\ % \ \tau \mapsto 1,\\
\mathbf{Re}(\textup{BPGL}) & \simeq \BP.
\end{align*}

%(\cite{HuKrizOrmsby, DuggerIsaksen})
The motivic dual Steenrod algebra over $\mathbb{C}$ is a Hopf algebra over $\pi_{*,*}(\F_p^{\textup{mot}}) = \F_p[\tau]$
which takes the following form at $p=2$ (see \cite[Sec. 12]{Voe03})
$$\A^\textup{mot} \cong \F_2[\tau][\tau_1, \ \tau_2, \cdots][\xi_1^2, \ \xi_2^2, \cdots]/\tau_i^2 = \tau \xi_i^2,$$
$$\mathbf{Re}(\tau)=1, \ \mathbf{Re}(\tau_i) = \xi_i, \ \mathbf{Re}(\xi_i^2) = \xi_i^2.$$
We then have the motivic Adams spectral sequence at every prime (\cite{DuggerIsaksen}) of the form
$$ \Ext_{\A^\textup{mot}}^{a,t,w}(\F_p[\tau], \F_p[\tau]) \cong E_2^{a,t,w} \Longrightarrow \pi_{t-a,w}S^{0,0}$$
$$d_r:E_r^{a,t,w} \rightarrow E_r^{a+r, t+r-1,w}$$
% Here $\A^{\textup{mot}}$ is motivic mod $p$ dual Steenrod algebra, $\pi_{*,*}(\F_p^{\textup{mot}}) = \F_p[\tau]$, and $\mathbf{Re}(\tau) = 1$.
The quotient map $S^{0,0} \rightarrow S^{0,0}/\tau$ induces a map of the motivic Adams spectral sequences.
$$\Ext_{\A^\textup{mot}}^{a,t,w}(\F_p[\tau],\, \F_p[\tau]) \longrightarrow \Ext_{\A^\textup{mot}}^{a,t,w}(\F_p[\tau],\, \F_p).$$

The following theorem is crucial in the computation of classical and motivic stable homotopy groups of spheres over $\mathbb{C}$.
\begin{thm}[{\cite[Theorem 1.17]{GWX}}] \label{MASS algNSS}
  There is an isomorphism of tri-graded spectral sequences between
  the motivic Adams spectral sequence for $S^{0,0}/\tau$ and
  the algebraic Novikov spectral sequence.
  \begin{displaymath}
    \xymatrix{
      \Ext_{{\A}^{\textup{mot}}}^{s+k, t'-k, \frac{t'}{2}}(\F_p[\tau],\, \F_p) \ar@{=>}[dd]|{\mathbf{Motivic \ Adams \ SS}}  \ar[rr]^{\cong} & &  \Ext_{\BP_*\BP/I}^{s,t'}(\F_p,\, I^{k}/I^{k+1}) \ar@{=>}[dd]|{\mathbf{Algebraic \ Novikov \ SS}} \\
      & & \\
      \pi_{t'-s, \frac{t'}{2}}(S^{0,0}/\tau)  \ar[rr]^-{\cong} & &   \Ext_{\BP_*\BP}^{s, t'}(\BP_*,\, \BP_*).
    }
  \end{displaymath}
\end{thm}

\begin{rmk}  
  Note that the $\Ext$-groups in the left column in Theorem~\ref{MASS algNSS} are only defined when $t'$ is even, meanwhile in the right column the $\Ext$-groups vanish for $t'$ odd for sparsity reasons.
\end{rmk}

As explained in \cite[Subsection~1.3]{GWX}, the motivic deformation and the naturality of the Adams spectral sequences give us a zig-zag diagram.
\begin{displaymath}
  \xymatrix{
\Ext_{\A}^{a,t}(\F_p,\, \F_p) \ar@{=>}[dd]|{\mathbf{Adams \ SS}}  & & 
  \Ext^{a,t,w}_{\A^\textup{mot}}(\F_p[\tau],\, \F_p[\tau]) \ar[ll]_-{\mathbf{Re}} \ar[rr] \ar@{=>}[dd]|{\mathbf{Motivic \ Adams \ SS}} & & 
   \Ext^{a,t,w}_{\A^\textup{mot}}(\F_p[\tau],\, \F_p)   \ar@{=>}[dd]|{\mathbf{Motivic \ Adams \ SS}} \\
 & & & & \\
 \pi_{t-a}S^0 & & \pi_{t-a,w}S^{0,0} \ar[ll]_-{\mathbf{Re}} \ar[rr] & & \pi_{t-a,w}S^{0,0}/\tau
}
\end{displaymath}
The right-hand horizontal maps are induced by the quotient map $S^{0,0} \rightarrow S^{0,0}/\tau$, and
the left-hand horizontal maps are given by the Betti realization functor.
The diagram of spectral sequences allow us to build up connections between the differentials in the classical Adams spectral sequence and the algebraic Novikov spectral sequence (by Theorem~\ref{MASS algNSS}) through the motivic world.

\begin{rmk} \label{tridegree}
At $p=2$, it is often useful to combine the isomorphisms in the Miller square and in Theorem~\ref{MASS algNSS} and obtain the following isomorphism between the $E_2$-pages of the Cartan-Eilenberg spectral sequence and the motivic Adams spectral sequence for $S^{0,0}/\tau$:
\[\xymatrix{
  \Ext^{s,t}_{\P}(\F_2, \Ext_{\mathcal{Q}}^k(\F_2, \F_2)) \ar[rr]^-{\cong} & & \Ext^{s+k, t, \frac{t-k}{2}}_{\A^{\textup{mot}}}(\F_2[\tau], \F_2).
}\]
\end{rmk}

We also have the motivic Adams-Novikov spectral sequence (\cite{HuKrizOrmsby, Isaksen}) of the form.
$$ \Ext_{\textup{BPGL}_{*,*}\textup{BPGL}}^{s,t',w}(\textup{BPGL}_{*,*},\, \textup{BPGL}_{*,*}) \cong E_2^{s,t',w} \Longrightarrow \pi_{t'-s,w}S^{0,0}$$
$$d_r:E_r^{s,t',w} \rightarrow E_r^{s+r, t'+r-1,w}$$

In \cite[Sections~6.1, 6.2]{Isaksen}, Isaksen proves the following rigidity theorem.
\begin{thm} \label{MANSS Isaksen}
After a re-grading, the motivic Adams-Novikov spectral sequence for $\pi_{*,*}S^{0,0}$ is isomorphic to a $\tau$-bockstein spectral sequence.
\end{thm}

Finally, for some of our later arguments we need a motivic version of the Cartan--Eilenberg spectral sequence.
Notably, this spectral sequence satisfies a rigidity theorem analogous to \Cref{MANSS Isaksen}.

\begin{cnstr}
  At $p=2$, we define 
  \begin{align*}
    \P^\textup{mot} & \coloneqq \F_2[\tau][\xi_1^2, \ \xi_2^2, \cdots] \cong \P[\tau], \\ 
    \mathcal{Q}^{\textup{mot}} & \coloneqq \A^{\textup{mot}}\otimes_{\P^{\textup{mot}}} \F_2[\tau] \cong \F_2[\tau][\tau_1, \ \tau_2, \cdots]
  \end{align*}
  The associated extension of Hopf algebroids
  $\P^\textup{mot} \rightarrow \A^\textup{mot} \rightarrow \mathcal{Q}^\textup{mot}$
  gives us a motivic version of the classical Cartan-Eilenberg spectral sequence.
  It has the form
  $$ \Ext^{s,t,w}_{\P^\textup{mot}}(\F_2[\tau],\, \Ext^k_{\mathcal{Q}^\textup{mot}}(\F_2[\tau],\, \F_2[\tau])) \cong E_2^{s,k,t,w} \Rightarrow \Ext^{s+k, t, w}_{\A^\textup{mot}}(\F_2[\tau],\, \F_2[\tau]), $$
$$d_r: E_r^{s,k,t,w} \to E_r^{s+r, k-r+1, t, w}.$$
  Note that in this spectral sequence all $d_r$-differentials look like motivic Adams $d_1$-differential in the tri-gradings.
\end{cnstr}

%% In Section~\ref{sec:mcess}, from an extension of Hopf algebroids $\P^{\textup{mot}} \rightarrow \A^{\textup{mot}} \rightarrow \mathcal{Q}^{\textup{mot}}$, we set up a motivic Cartan-Eilenberg spectral sequence that converges to the motivic Adams $E_2$-page, and prove the following rigidity theorem in analogue of Theorem~\ref{MANSS Isaksen}.

\begin{thm}[Rigidity Theorem] \label{MCESS tau Bockstein}
At $p=2$, after a re-grading, the motivic Cartan-Eilenberg spectral sequence for $ \Ext_{{\A}^{\textup{mot}}}^{*, *, *}(\F_2[\tau], \F_2[\tau])$ is isomorphic to a $\tau$-Bockstein spectral sequence.
\end{thm}

\begin{proof}
We have 
\begin{align*}
\Ext^{s,t,w}_{\P^{\textup{mot}}}(\F_2[\tau], \Ext^k_{\mathcal{Q}^\textup{mot}}(\F_2[\tau], \F_2[\tau])) & \cong \Ext^{s,t,w}_{\P^\textup{mot}}(\F_2[\tau], \Ext^k_{\mathcal{Q}}(\F_2, \F_2))[\tau])  \\ 
& \cong \Ext_{\P}^{s,t}(\F_2, \Ext^k_{\mathcal{Q}}(\F_2, \F_2))[\tau].
\end{align*}
Here $\tau$ has $(s,k,t,w)$-degrees $(0,0,0, -1)$, and every element in $\Ext_{\P}^{s,t}(\F_2, \Ext^k_{\mathcal{Q}}(\F_2, \F_2))$ satisfies $w = \frac{t-k}{2}$ by Remark~\ref{tridegree}. Note that by sparseness this group is nonzero only if $t-k$ is even.

Consider $S^{0,0}/\tau$, we have an isomorphism of the $E_2$-page of its MCESS 
\begin{align*}
\Ext^{s,t,w}_{\P^{\textup{mot}}}(\F_2[\tau], \Ext^k_{\mathcal{Q}^\textup{mot}}(\F_2[\tau], \F_2)) & \cong \Ext^{s,t,w}_{\P^\textup{mot}}(\F_2[\tau], \Ext^k_{\mathcal{Q}}(\F_2, \F_2)))  \\ 
& \cong \Ext_{\P}^{s,t}(\F_2, \Ext^k_{\mathcal{Q}}(\F_2, \F_2)).
\end{align*}
Since all differentials in MCESS preserve the $w$ and $t$-degrees, it must preserve the $k$-degree in the case for $S^{0,0}/\tau$, so its MCESS collapses at the $E_2$-page.

Back to $S^{0,0}$, again due to degree reasons, all $d_r$-differentials have the form
$$d_{2n+1} x = \tau^n y, \ \ x, y \in \Ext_{\P}^{s,t}(\F_2, \Ext^k_{\mathcal{Q}}(\F_2, \F_2)),$$
and $d_{2n} = 0$. One can check that there are no $\tau$-extensions in MCESS. This is precisely saying that after a regrading, the MCESS for $\Ext^{*, *, *}_{\A^\textup{mot}}(\F_2[\tau], \F_2[\tau])$ is isomorphic to a $\tau$-Bockstein spectral sequence and completes the proof.
\end{proof}

\begin{rmk}
The important consequence of Theorem~\ref{MANSS Isaksen} is that, the differentials in the classical Adams-Novikov spectral sequence completely determine the differentials in the motivic Adams-Novikov spectral sequence, and vice versa. Similarly, our Theorem~\ref{MCESS tau Bockstein} tells us that the differentials in the classical and motivic Cartan-Eilenberg spectral sequences determine each other.
\end{rmk}

\begin{cor} \label{cor:no entering}
  Suppose there are no non-trivial Cartan--Eilenberg differentials entering tridegrees
  $(s,k,t) = (s,*,t)$.
  Let $\{\overline{x}_i\}_{i \in I}$ be a collection of generators of the permanent cycles on the Cartan--Eilenberg $E_2$-page in tridegree $(s,k,t) = (s,*,t)$.
  Then, we can lift each $\overline{x}_i$ to a class $x_i$ on the $E_2$-page of the motivic Adams spectral sequence for $S^{0,0}$ and any choice of lifts $\{x_i\}_{i \in I}$ provides a basis for $E_2^{s,*,t}$ as a free $\F_2[\tau]$-module.
\end{cor}

%This fact is very useful in our main arguments in Section~\ref{sec:proof}.
